% -*- mode:LaTex; mode:visual-line; mode:flyspell; fill-column:75-*-

\chapter{Conclusions} \label{chapConclusions}

% TODO: Briefly restate the problem and what was accomplished.
% 1--2 paragraphs summarizing the thesis at a high level.

\section{Summary of Contributions} \label{secContribSummary}

% TODO: Revisit the contributions listed in Section~\ref{secContributions}
% and describe what was actually accomplished for each.

\begin{enumerate}
  \item \textbf{Multi-modal orchard dataset.}
  % TODO: Summarize dataset size, modalities, and key properties.

  \item \textbf{2D fire blight detection pipeline.}
  % TODO: Summarize best detection performance (mAP, modality comparison).

  \item \textbf{Semantic OctoMap for disease mapping.}
  % TODO: Summarize the representation and its fidelity in simulation/field.

  \item \textbf{Semantic Next-Best-View planner.}
  % TODO: Summarize improvement over baselines (simulation and field).

  \item \textbf{Full-system field validation.}
  % TODO: Summarize field deployment results and what was demonstrated.
\end{enumerate}

\section{Limitations} \label{secLimitations}

% TODO: Discuss limitations of the current system honestly.
% - Detection accuracy limitations (small lesions, early infection).
% - Map accuracy vs. localization error.
% - Real-time performance constraints.
% - Generalization to other orchards, tree varieties, or disease severities.
% - Assumption of static scene during a mission.

\section{Future Work} \label{secFutureWork}

% TODO: Suggest directions for future research building on this thesis.

\subsection{Improved Detection Models} \label{secFutureDetection}

% TODO: e.g., larger or more diverse training data; few-shot adaptation
% to new disease stages or varieties; real-time NIR + RGB fusion models.

\subsection{Richer Map Representations} \label{secFutureMapping}

% TODO: e.g., Gaussian splatting as the live map (not just for GT);
% neural radiance fields for appearance modeling; temporal maps
% tracking disease progression across visits.

\subsection{Long-Horizon Planning and Multi-Robot Systems} \label{secFuturePlanning}

% TODO: e.g., multi-day mission planning; coordinating a fleet of robots
% across a full orchard; integration with orchard management software
% for automated pruning prescription.

\subsection{Closed-Loop Disease Management} \label{secFutureManagement}

% TODO: Moving beyond detection toward action: automated targeted pruning,
% integration with spray systems, feedback to orchard managers.

\section{Closing Remarks} \label{secClosingRemarks}

% TODO: Final paragraph — broader impact statement.
% What does this work contribute to agricultural robotics and
% plant disease management? What vision does it open up?
